% Book-local P12.6 normalized-link functoriality for finite curve maps.
\begin{proposition}[Finite curve maps on normalized branch links]
\label{prop:p126-finite-reduced-curve-link}
Let
\[
 \varphi:(R,\mathfrak m)\longrightarrow(S,\mathfrak n)
\]
be a finite local homomorphism between one-dimensional reduced Nagata
noetherian local rings, and let \(\mathfrak a\subset R\) be
\(\mathfrak m\)-primary.  Then \(\mathfrak aS\) is
\(\mathfrak n\)-primary and restriction defines a continuous map
\[
 \varphi^*: \mathcal V_{\mathfrak aS}(S)\longrightarrow
             \mathcal V_{\mathfrak a}(R),
 \qquad w\longmapsto w\circ\varphi.
 \tag{\(\ddagger\)}\label{eq:p126-finite-curve-link-map}
\]

More precisely, let \(\mathfrak q'\in\operatorname{Br}(S,\mathfrak n)\).
It lies over a unique branch
\(\mathfrak q\in\operatorname{Br}(R,\mathfrak m)\), and the induced local
map of discrete valuation rings has a ramification index
\(e_{\mathfrak q'/\mathfrak q}\geq1\).  With the notation of
\Cref{thm:p121-reduced-nagata-curve-link},
\[
 \begin{aligned}
 w_{\mathfrak q'}(\varphi(f))
   &=e_{\mathfrak q'/\mathfrak q}\,w_{\mathfrak q}(f)
       &&(f\in R),\\
 d_{\mathfrak q'}(\mathfrak aS)
   &=e_{\mathfrak q'/\mathfrak q}\,
       d_{\mathfrak q}(\mathfrak a),\\
 \varphi^*\!\left(
   \frac{w_{\mathfrak q'}}{d_{\mathfrak q'}(\mathfrak aS)}
                 \right)
   &=\frac{w_{\mathfrak q}}{d_{\mathfrak q}(\mathfrak a)}.
 \end{aligned}
 \tag{\(\heartsuit\)}\label{eq:p126-ramification-cancellation}
\]
Here the first equality uses the convention
\(e\cdot(+\infty)=+\infty\).  These maps commute with composition of
finite local homomorphisms and with the change-of-ideal maps
\eqref{eq:p121-change-normalizing-ideal}.
\end{proposition}

\begin{proof}
The quotient \(S/\mathfrak aS\) is finite over the artinian ring
\(R/\mathfrak a\), hence is artinian.  It is a nonzero local ring, since
\(\mathfrak aS\subseteq\mathfrak mS\subseteq\mathfrak n\).  Its unique
prime is therefore \(\mathfrak n/\mathfrak aS\), proving that
\(\mathfrak aS\) is \(\mathfrak n\)-primary.  For every nonnegative
semivaluation \(w\) on \(S\),
\[
 w(\mathfrak aS)=\min_{f\in\mathfrak a}w(\varphi(f)),
 \tag{\(\clubsuit\)}\label{eq:p126-extended-ideal-value}
\]
because coefficients from \(S\) have nonnegative value and the coefficient
one occurs.  Hence restriction preserves the displayed normalization and
is continuous for pointwise convergence.  Since \(\varphi\) is local,
positivity on \(\mathfrak n\) restricts to positivity on \(\mathfrak m\), so
the restricted semivaluation remains centred at the closed point.

Let \(\mathfrak p'\) be the minimal prime of \(S\) corresponding to
\(\mathfrak q'\), and put
\(\mathfrak p=\varphi^{-1}(\mathfrak p')\).  The induced map
\[
             R/\mathfrak p\lhook\joinrel\longrightarrow
             S/\mathfrak p'
\]
is finite and integral.  Since the target is a one-dimensional domain,
dimension invariance for integral injective extensions
\cite[Tag~00OK]{StacksProject} shows that
\(\dim(R/\mathfrak p)=1\).  Thus \(\mathfrak p\) is a minimal prime of
\(R\).

The inclusion of the component fraction fields carries the normalization
of \(R/\mathfrak p\) into that of \(S/\mathfrak p'\).  Contracting the
centre \(\mathfrak q'\) and localizing therefore gives an injective local
map
\[
                       D_{\mathfrak q}\longrightarrow D_{\mathfrak q'}
\]
of discrete valuation rings.  By the definition of ramification index
\cite[Tag~09E4]{StacksProject}, a uniformizer of the first ring is a unit
times the \(e_{\mathfrak q'/\mathfrak q}\)-th power of a uniformizer of
the second.  This proves the first line of
\eqref{eq:p126-ramification-cancellation}.  Taking the minimum over
\(\mathfrak a\) and using \eqref{eq:p126-extended-ideal-value} proves the
second; dividing cancels the ramification index and proves the third.

Restriction of functions is associative.  Equivalently, the branchwise
ramification indices multiply in a tower by
\cite[Tag~0BRL]{StacksProject}.  This proves compatibility with composition.
The branchwise description of
\eqref{eq:p121-change-normalizing-ideal} proves compatibility with changing
the normalizing ideal.
\end{proof}
